\documentclass[11pt,a4paper]{article}

% =================================================
% PACKAGES
% =================================================

\usepackage[margin=2.5cm]{geometry}
\usepackage{booktabs}
\usepackage{longtable}
\usepackage{hyperref}
\usepackage{xcolor}
\usepackage{enumitem}
\usepackage{fancyhdr}

% =================================================
% PAGE SETUP
% =================================================

\pagestyle{fancy}
\fancyhf{}

\fancyhead[L]{Cybersecurity Assessment Report}
\fancyhead[R]{MLDA Swarm}

\fancyfoot[C]{\thepage}

\setlength{\parindent}{0pt}
\setlength{\parskip}{6pt}

% =================================================
% DOCUMENT
% =================================================

\begin{document}

% =================================================
% TITLE PAGE
% =================================================

\begin{titlepage}

    \centering

    \vspace*{3cm}

    {\Huge\bfseries Cybersecurity Assessment Report\par}

    \vspace{1cm}

    {\Large MLDA Adaptive Cybersecurity Swarm\par}

    \vfill

    \begin{tabular}{ll}

        \textbf{Run ID:}
        &
        {{ report.run_id | latex }}
        \\

        \textbf{Target(s):}
        &
        {{ report.target_ids | join(", ") | latex }}
        \\

        \textbf{Generated At:}
        &
        {{ report.generated_at | latex }}
        \\

    \end{tabular}

    \vfill

    {\large Confidential Security Assessment}

\end{titlepage}


% =================================================
% TABLE OF CONTENTS
% =================================================

\tableofcontents

\newpage


% =================================================
% ASSESSMENT INFORMATION
% =================================================

\section{Assessment Information}

\begin{itemize}[leftmargin=*]

    \item \textbf{Run ID:}
    {{ report.run_id | latex }}

    \item \textbf{Target(s):}
    {{ report.target_ids | join(", ") | latex }}

    \item \textbf{Generated At:}
    {{ report.generated_at | latex }}

\end{itemize}


% =================================================
% EXECUTIVE SUMMARY
% =================================================

\section{Executive Summary}

{% if report.executive_summary %}

{{ report.executive_summary | latex }}

{% else %}

Executive summary has not yet been generated.

{% endif %}


% =================================================
% SEVERITY SUMMARY
% =================================================

\section{Severity Summary}

\begin{table}[h]
    \centering

    \begin{tabular}{lc}

        \toprule

        \textbf{Severity}
        &
        \textbf{Number of Findings}
        \\

        \midrule

        Critical
        &
        {{ report.severity_summary.critical }}
        \\

        High
        &
        {{ report.severity_summary.high }}
        \\

        Medium
        &
        {{ report.severity_summary.medium }}
        \\

        Low
        &
        {{ report.severity_summary.low }}
        \\

        Informational
        &
        {{ report.severity_summary.info }}
        \\

        \midrule

        \textbf{Total}
        &
        \textbf{ {{ report.severity_summary.total }} }
        \\

        \bottomrule

    \end{tabular}

    \caption{Summary of identified security findings}

\end{table}


% =================================================
% DETAILED FINDINGS
% =================================================

\newpage

\section{Detailed Findings}


{% if report.findings %}


{% for finding in report.findings %}

% -------------------------------------------------
% Finding
% -------------------------------------------------

\subsection{ {{ finding.title | latex }} }

\begin{itemize}[leftmargin=*]

    \item \textbf{Finding ID:}
    {{ finding.finding_id | latex }}

    \item \textbf{Severity:}
    {{ finding.severity | capitalize | latex }}

    \item \textbf{CVSS Score:}

    {% if finding.cvss_score is not none %}

    {{ finding.cvss_score }}

    {% else %}

    Not scored

    {% endif %}

    \item \textbf{Affected Component:}
    {{ finding.affected_component | latex }}

    \item \textbf{Detected By:}
    {{ finding.agent_name | latex }}


    {% if finding.tool_name %}

    \item \textbf{Tool:}
    {{ finding.tool_name | latex }}

    {% endif %}

\end{itemize}


% -------------------------------------------------
% Description
% -------------------------------------------------

\subsubsection*{Description}

{{ finding.description | latex }}


% -------------------------------------------------
% Evidence
% -------------------------------------------------

\subsubsection*{Evidence}

{% if finding.evidence %}

\begin{quote}

{{ finding.evidence | latex }}

\end{quote}

{% else %}

No additional evidence provided.

{% endif %}


% -------------------------------------------------
% Recommendation
% -------------------------------------------------

\subsubsection*{Recommendation}

{% if finding.recommendation %}

{{ finding.recommendation | latex }}

{% else %}

No remediation recommendation provided.

{% endif %}


% -------------------------------------------------
% References
% -------------------------------------------------

{% if finding.references %}

\subsubsection*{References}

\begin{itemize}

{% for reference in finding.references %}

    \item {{ reference | latex }}

{% endfor %}

\end{itemize}

{% endif %}


\vspace{0.5cm}

\hrule

\vspace{0.5cm}


{% endfor %}


{% else %}

No security findings were identified during this assessment.

{% endif %}


% =================================================
% GENERAL RECOMMENDATIONS
% =================================================

\newpage

\section{General Recommendations}

The following remediation approach is recommended based on the
findings identified during the assessment:

\begin{enumerate}

    \item Address Critical severity findings immediately.

    \item Prioritise High severity vulnerabilities after Critical
    findings have been addressed.

    \item Review and remediate Medium severity findings according to
    the affected system's risk profile.

    \item Patch outdated software and exposed services.

    \item Review authentication and access-control mechanisms.

    \item Rotate credentials if credential exposure has been detected.

    \item Conduct a follow-up security assessment after remediation.

\end{enumerate}


% =================================================
% REPORT INFORMATION
% =================================================

\section{Report Information}

This report was generated by the MLDA Adaptive Cybersecurity Swarm.

The assessment is associated with the following run identifier:

\begin{quote}

\textbf{Run ID:}
{{ report.run_id | latex }}

\end{quote}

Sensitive information identified during assessment processing may be
redacted from the exported report.

Where redaction has been applied, sensitive values may appear as:

\begin{quote}

\texttt{[REDACTED]}

\end{quote}

Operational logs, prompts, tool-call details, latency measurements,
token usage and other internal execution metadata are maintained
separately from this assessment report.


% =================================================
% END DOCUMENT
% =================================================

\end{document}